\section{Objetivos del Sistema}
\label{sec:SystemObjectives}

A partir de los objetivos de negocio anteriores, el sistema desarrollado persigue la construcción de una \textbf{plataforma web unificada} que permite al distribuidor, a su red comercial y a los clientes profesionales operar sobre una misma base de información. De esta manera, se evita la fragmentación entre herramientas aisladas de catálogo, seguimiento comercial y gestión de pedidos.

De forma más concreta, los objetivos generales del sistema son los siguientes:

\begin{itemize}
    \item \textbf{Centralizar en una única aplicación} la gestión de usuarios, clientes, catálogo, visitas, rutas, pedidos y seguimiento de cobros.
    \item \textbf{Garantizar un acceso diferenciado por roles}, de manera que cada perfil utilice una interfaz y un conjunto de operaciones acordes a sus responsabilidades reales dentro del proceso comercial.
    \item \textbf{Facilitar al distribuidor la administración operativa} del sistema, incluyendo la gestión de usuarios, la auditoría de acciones relevantes, la consulta global de pedidos, la configuración de avisos automáticos y las bases técnicas necesarias para futuras integraciones empresariales.
    \item \textbf{Proporcionar al comercial una herramienta móvil y operativa} para planificar su jornada, consultar la ficha de sus clientes, preparar pedidos, organizar repartos, imprimir etiquetas y registrar el resultado de sus visitas.
    \item \textbf{Permitir al cliente profesional consultar información clave} como el catálogo, las cartas de color, los pedidos, el historial y el estado logístico de sus entregas cuando exista un reparto real planificado.
    \item \textbf{Mantener trazabilidad sobre los procesos críticos} del sistema, especialmente en autenticación, gestión administrativa de usuarios, ciclo de vida del pedido y seguimiento de cobros.
    \item \textbf{Integrar servicios externos de soporte} cuando aporten valor claro al flujo funcional, como almacenamiento de imágenes, geocodificación, cálculo de rutas, generación de QR o soporte de instalación como \textit{aplicación web progresiva}.
    \item \textbf{Asegurar una evolución incremental} de la aplicación, manteniendo una base modular que permita ampliar funcionalidades futuras sin rehacer los módulos ya estabilizados.
\end{itemize}
